\documentclass[uplatex,a4paper,12pt]{jsarticle}
\usepackage{amsmath,amssymb}
\usepackage{enumitem}

\begin{document}

\title{機械学習基礎 確認テスト 解答}
\author{}
\date{}
\maketitle

\section*{配点}
全10問，各10点，合計100点．60点以上を合格とする．

\section*{解答}

\begin{enumerate}

\item
\[
\boxed{(b)}
\]
機械学習は，データ集合からパターンや判断ルールを自動で獲得する手法である．

\item
\[
\boxed{(b)}
\]
教師あり学習では，入力と出力のペア $(x,y)$ を用いる．

\item
\[
\boxed{(c)}
\]
株価予測は連続値を予測するため，回帰問題である．

\item
\[
\boxed{(b)}
\]
分類問題は，入力から離散的なクラスを予測する問題である．

\item
\[
\boxed{(b)}
\]
単回帰モデル
\[
y=\alpha+\beta x
\]
において，$\alpha$ はバイアス，$\beta$ は重みに相当する．

\item
\[
\boxed{(a)}
\]
未知のデータに対する予測性能を汎化性能という．

\item
\[
\boxed{(a)}
\]
過学習とは，学習データに過剰に適合し，未知データで性能が低下する現象である．

\item
正解率は
\[
Accuracy=
\frac{TP+TN}{TP+TN+FP+FN}
\]
である．

したがって，
\[
Accuracy=
\frac{30+50}{30+50+10+10}
=
\frac{80}{100}
=
0.8
\]

よって，
\[
\boxed{0.8}
\]
または
\[
\boxed{80\%}
\]

\item
\[
\boxed{(a)}
\]
ソフトマックス関数は，各クラスに属する確率を出力する．

\item
\[
\boxed{(a)}
\]
学習率 $\eta$ は，勾配降下法におけるパラメータの更新量を調整する値である．

\end{enumerate}

\end{document}